\documentclass[11pt,spanish]{article}

% =========================================================
% PLANTILLA BASE PARA INFORMES ACADÉMICOS / TÉCNICOS
% =========================================================
% Instrucciones rápidas:
% 1. Edite los datos de la sección "DATOS DEL DOCUMENTO".
% 3. Use las tablas de ejemplo como base para nuevas tablas.
% 4. Para tablas anchas, use el entorno widepage incluido abajo.
% 5. Las imágenes se cargan con \SafeImage, que no falla si el archivo no existe.

% --- Idioma y codificación ---
\usepackage[spanish,es-tabla]{babel}
\usepackage[T1]{fontenc}
\usepackage[utf8]{inputenc}

% --- Márgenes / papel ---
\usepackage{geometry}
\geometry{a4paper, margin=2.7cm}

% --- Matemáticas ---
\usepackage{amsmath}

% --- Tablas, enlaces y elementos visuales ---
\usepackage{array}
\usepackage{tabularx}
\usepackage{booktabs}
\usepackage{longtable}
\usepackage{graphicx}
\usepackage{float}
\usepackage{xcolor}
\usepackage{tikz}
\usetikzlibrary{positioning,shapes.geometric,arrows.meta,calc}
\usepackage{enumitem}
\usepackage{ragged2e}
\usepackage[hidelinks]{hyperref}

% --- Encabezado y pie de página ---
\usepackage[myheadings]{fancyhdr}
\usepackage{lastpage}
\setlength{\headheight}{30pt}
\setlength{\footskip}{30pt}

% --- Interlineado ---
\linespread{1.2}
\sloppy

% =========================================================
% DATOS DEL DOCUMENTO
% =========================================================
\newcommand{\Institucion}{Instituto Tecnológico de Costa Rica}
\newcommand{\Campus}{Campus Tecnológico Central Cartago}
\newcommand{\DocumentoTitulo}{Entregable 3 -- Registro de Stakeholders}
\newcommand{\DocumentoSubtitulo}{}
\newcommand{\Curso}{Administración de proyectos}
\newcommand{\Profesor}{Alicia Marcela Salazar Hernández}
\newcommand{\FechaDocumento}{7 de junio de 2026}
\newcommand{\PeriodoAcademico}{I Semestre}
\newcommand{\AnioDocumento}{2026}

\newcommand{\EstudiantesPortada}{%
    Mauricio González Prendas: 2024143009\\
    Isaac Jiménez Blanco: 2024202804\\
    Tamara Robles Camacho: 2024099342%
}

% Datos que aparecen en el encabezado.
\newcommand{\DocHeaderLeft}{}
\newcommand{\DocHeaderRight}{Registro de Stakeholders}
\newcommand{\DocFooterRight}{Página~\thepage\ de~\pageref{LastPage}}

% Metadatos PDF.
\title{\DocumentoTitulo}
\author{\EstudiantesPortada}
\date{\FechaDocumento}

% =========================================================
% CONFIGURACIÓN DE TABLAS Y UTILIDADES
% =========================================================
\newcolumntype{Y}{>{\RaggedRight\arraybackslash}X}
\renewcommand{\arraystretch}{1.22}
\setlength{\LTpre}{0.5em}
\setlength{\LTpost}{0.5em}

% Imagen segura: si el archivo no existe, muestra un recuadro de marcador.
\newcommand{\SafeImage}[2][0.92\textwidth]{%
    \IfFileExists{#2}{%
        \includegraphics[width=#1]{#2}%
    }{%
        \fbox{%
            \begin{minipage}[c][4cm][c]{#1}
            \centering
            Imagen pendiente:\\
            \texttt{#2}
            \end{minipage}%
        }%
    }%
}

% Figura reutilizable.
\newcommand{\EvidenceFigure}[3]{%
    \begin{figure}[H]
    \centering
    \SafeImage{#1}
    \caption{#2}
    \label{#3}
    \end{figure}%
}

% =========================================================
% ENCABEZADO Y PIE
% =========================================================
\pagestyle{fancy}
\fancyhf{}
\fancyhead[L]{\DocHeaderLeft}
\fancyhead[R]{\DocHeaderRight}
\renewcommand{\headrulewidth}{0.4pt}
\renewcommand{\footrulewidth}{0.4pt}
\fancyfoot[R]{\DocFooterRight}

% Estilo equivalente para páginas horizontales.
\fancypagestyle{widefancy}{%
    \fancyhf{}%
    \fancyhead[L]{\DocHeaderLeft}%
    \fancyhead[R]{\DocHeaderRight}%
    \renewcommand{\headrulewidth}{0.4pt}%
    \renewcommand{\footrulewidth}{0.4pt}%
    \fancyfoot[R]{\DocFooterRight}%
}

% =========================================================
% PÁGINAS HORIZONTALES PARA TABLAS ANCHAS
% =========================================================
\makeatletter
\newcommand{\SyncPageBuilderDimensions}{%
    \global\vsize=\textheight
    \global\@colroom=\textheight
    \global\@colht=\textheight
}
\makeatother

\newenvironment{widepage}{%
    \clearpage
    \hypersetup{pageanchor=false}%
    \paperwidth=297mm
    \paperheight=210mm
    \pdfpagewidth=297mm
    \pdfpageheight=210mm
    \newgeometry{left=2.7cm,right=2.7cm,top=2.7cm,bottom=2.7cm}%
    \setlength{\textwidth}{243mm}%
    \setlength{\textheight}{156mm}%
    \setlength{\linewidth}{\textwidth}%
    \setlength{\columnwidth}{\textwidth}%
    \hsize=\textwidth
    \setlength{\headwidth}{\textwidth}%
    \SyncPageBuilderDimensions
    \pagestyle{widefancy}%
}{%
    \clearpage
    \paperwidth=210mm
    \paperheight=297mm
    \pdfpagewidth=210mm
    \pdfpageheight=297mm
    \restoregeometry
    \setlength{\textwidth}{156mm}%
    \setlength{\textheight}{243mm}%
    \setlength{\linewidth}{\textwidth}%
    \setlength{\columnwidth}{\textwidth}%
    \hsize=\textwidth
    \setlength{\headwidth}{\textwidth}%
    \SyncPageBuilderDimensions
    \hypersetup{pageanchor=true}%
    \pagestyle{fancy}%
}

% =========================================================
% PORTADA
% =========================================================
\newcommand{\MakeCover}{%
\begin{titlepage}
\thispagestyle{empty}
\centering

{\bfseries \huge \Institucion \par}
\vspace{0.25cm}
{\LARGE \Campus \par}

\vfill

{\huge Registro de Stakeholders \par}

\vfill

{\bfseries \LARGE Profesora: \par}
{\Large \Profesor \par}

\vfill

{\bfseries \LARGE Estudiantes: \par}
{\Large \EstudiantesPortada \par}

\vfill

{\bfseries\LARGE Fecha: \par}
{\Large \FechaDocumento \par}

\vfill

{\LARGE \PeriodoAcademico \par}
{\LARGE \AnioDocumento \par}

\end{titlepage}%
}

% =========================================================
% DOCUMENTO
% =========================================================
\begin{document}
\hypersetup{pageanchor=false}

% =========================
% Portada
% =========================
\MakeCover

% =========================
% Índice
% =========================
\pagenumbering{roman}
\pagestyle{empty}
\tableofcontents
\clearpage
\pagenumbering{arabic}
\hypersetup{pageanchor=true}
\pagestyle{fancy}

% =========================
% Contenido
% =========================

\section{Introducción}

En este documento corresponde al Registro de Stakeholders (partes interesadas o interesados) del proyecto Plataforma Universitaria Integrada, desarrollado para la Universidad Tecnológica La Mejor. Esta información forma parte del grupo de procesos de Inicio, según el marco metodológico del PMI, y su propósito es identificar, clasificar y definir a las personas que formaran parte de este proyecto y que pueden afectar o verse afectados por el desarrollo y los resultados del mismo.


La identificación temprana de los interesados es fundamental para garantizar una buena comunicación, gestionar expectativas y asegurar el apoyo durante el ciclo de vida del proyecto. Los stakeholdes se dividirán en dos partes, los internos (miembros del equipo de desarrollo y administración) y los externos (futuros actores de la plataforma), para cada uno se documenta su cargo, nivel de interés, influencia y poder, así como la estrategia más adecuada para mantener su participación con los objetivos del proyecto.


La estácala de referencia para comprender mejor la tabla es: Influencia y Poder se valoran del 1 al 5 (1= muy bajo, 5= muy alto)

\section{Stakeholders internos del proyecto}

En la siguiente tabla se muestran los stakeholders Internos del proyecto:

\begin{widepage}
{\small
\setlength{\tabcolsep}{3pt}
\begin{longtable}{p{3.8cm} p{3.8cm} p{5.8cm} p{1.9cm} p{1.4cm} p{5.8cm}}
\toprule
\textbf{Nombre} & \textbf{Cargo} & \textbf{Interés} & \textbf{Influencia} & \textbf{Poder} & \textbf{Estrategia de involucramiento} \\
\midrule
\endfirsthead
\toprule
\textbf{Nombre} & \textbf{Cargo} & \textbf{Interés} & \textbf{Influencia} & \textbf{Poder} & \textbf{Estrategia de involucramiento} \\
\midrule
\endhead
Tamara Robles Camacho & Project Manager & Garantizar que el proyecto se ejecute dentro del alcance, tiempo y costo definidos & 5 & 5 & Liderar reuniones semanales de seguimiento, tomar decisiones sobre cambios y gestionar comunicaciones con el patrocinador \\
Mauricio González Prendas & Desarrollador Frontend 1 \& QA & Desarrollar un front-end funcional, navegable y de calidad que cumpla los criterios de aceptación & 5 & 4 & Participar en la definición del alcance técnico, revisiones de código y pruebas de calidad antes de cada entrega \\
Carlos Sainz Arce & Desarrollador Frontend 2 \& QA & Desarrollar un front-end funcional, navegable y de calidad que cumpla los criterios de aceptación & 5 & 4 & Participar en la definición del alcance técnico, revisiones de código y pruebas de calidad antes de cada entrega junto con el desarrollador 1 \\
Isaac Jiménez Blanco & Analista de Negocio \& Documentador 1 & Asegurar que todos los entregables documentales estén completos, coherentes y alineados con los estándares PMI & 5 & 4 & Responsable de redactar y revisar entregables, mantener actualizado el registro de riesgos y la documentación del proyecto \\
Jorge Abarca Quiros & Analista de Negocio \& Documentador 2 & Asegurar, junto con el analista y documentador 1, que todos los entregables documentales estén completos, coherentes y alineados con los estándares PMI & 5 & 4 & Redacta y analiza los entregables, manteniendo actualizado la documentación. \\
\bottomrule
\end{longtable}
}
\end{widepage}
\section{Stakeholders externos}

En la siguiente tabla se mencionan las partes externas interesadas para este proyecto:

\begin{widepage}
{\small
\setlength{\tabcolsep}{3pt}
\begin{longtable}{p{3.8cm} p{4.3cm} p{5.8cm} p{1.9cm} p{1.4cm} p{5.4cm}}
\toprule
\textbf{Nombre} & \textbf{Cargo} & \textbf{Interés} & \textbf{Influencia} & \textbf{Poder} & \textbf{Estrategia de involucramiento} \\
\midrule
\endfirsthead
\toprule
\textbf{Nombre} & \textbf{Cargo} & \textbf{Interés} & \textbf{Influencia} & \textbf{Poder} & \textbf{Estrategia de involucramiento} \\
\midrule
\endhead
Dr. Roberto Mora Fallas & Rector – Patrocinador del Proyecto & Modernizar la institución, mejorar la imagen universitaria y optimizar los procesos internos & 4 & 5 & Presentarle el Project Charter para aprobación, informes ejecutivos mensuales y acta de cierre \\
Lic. Patricia Vargas Soto & Directora Administrativa & Reducir la carga operativa del personal y eliminar duplicidad de datos entre departamentos & 4 & 4 & Reuniones de validación de requerimientos administrativos y revisión del módulo de gestión \\
Ing. Carlos Brenes Mora & Director de TI & Garantizar que la solución sea técnicamente viable, segura y mantenible a futuro & 3 & 4 & Consultas técnicas en la fase de planificación y revisión de la arquitectura del frontend \\
Prof. Andrea Solís Zamora & Representante Docente & Contar con herramientas digitales que simplifiquen el registro de calificaciones y asistencia & 4 & 3 & Participar en sesiones de validación del portal docente y pruebas de usuario \\
Estudiantes (colectivo) & Usuarios Finales – Portal Estudiantil & Acceder a matrícula, calificaciones, horarios y pagos de forma rápida y en línea & 5 & 2 & Encuestas de satisfacción y pruebas de usabilidad del portal estudiantil \\
Depto. Financiero & Área Financiera & Digitalizar los procesos de cobro de matrícula y generación de reportes financieros & 4 & 3 & Validación de requerimientos del módulo financiero y revisión de casos de uso así como las pruebas de usuario para ese modulo. \\
Proveedor de Hosting & Proveedor Externo & Ofrecer infraestructura confiable para el despliegue del sistema & 2 & 2 & Contrato de servicio con SLA definido, revisión periódica de disponibilidad \\
\bottomrule
\end{longtable}
}
\end{widepage}

\end{document}

